\section{Conclusion}
\label{sec:conclusion}

This work formulates two-view 3D pose fusion after monocular estimation. A6
combines pelvis-centered canonicalization, a transparent quality-weighted base,
a view-swap-invariant bounded temporal residual, and self-supervised recovery,
skeletal, temporal, rotation, and complete-cycle constraints without reading
triangulated targets during training or selection.

The evidence supports a narrower conclusion than positional superiority.
Canonicalization accounts for most improvement against the private same-video
pseudo-reference. On 14 held-out people, A6 is statistically
indistinguishable from arithmetic averaging and A5. It remains the method
mainline as the complete auditable constrained objective. On limited Unity
native 3D, A6 does not improve zero-shot direct-3D transfer, while calibrated
2D triangulation is substantially more accurate. This negative result marks the
boundary of uncalibrated post-estimation fusion rather than weakening the
importance of geometry.

OOF A6 features support a person-disjoint application over 928 cycles, but the
elderly--student speed and jerk coefficients are materially pose-source- and
estimand-sensitive, and no cohort-specific repetition trend survives
correction. The cohort analysis is therefore exploratory and observational.
The method is best interpreted as an auditable motion-representation layer;
diverse independent capture, repeated seeds, measured participant covariates,
and completed governance documentation are required before absolute,
biomechanical, causal, or clinical claims.
